\section{Prosedur, Pemanggilan, dan Pengembalian}
\label{sec:prosedur-pemanggilan}

Program yang ditulis sebagai satu deretan instruksi panjang sulit
dibaca, diuji, dan diperbaiki. Ketika tugas yang sama harus
dikerjakan pada bagian program yang berbeda, kodenya harus disalin
berkali-kali, sehingga setiap kekeliruan juga harus diperbaiki
berkali-kali. Solusinya adalah \term{prosedur}.

\subsection{Mengapa Prosedur Diperlukan}

Ada tiga alasan utama program dipecah menjadi prosedur.

\begin{enumerate}
    \item \term{Modularitas.} Setiap prosedur menangani satu tugas
          yang batasnya jelas, sehingga program besar menjadi
          kumpulan bagian kecil yang dapat diuji dan diperbaiki
          secara terpisah.

    \item \term{Penggunaan ulang} (\eng{reuse}). Kode yang sering
          dipakai cukup ditulis sekali. Prosedur pencetak bilangan
          heksadesimal, misalnya, dapat dipanggil dari puluhan tempat
          tanpa perlu disalin.

    \item \term{Abstraksi.} Pemanggil cukup mengetahui antarmuka
          prosedur: nama, parameter, dan nilai baliknya. Rincian di
          dalam prosedur dapat diganti tanpa mengubah pemanggilnya.
\end{enumerate}

Pembahasan Bab~\ref{chap:kontrol} menunjukkan bahwa \asm{JMP} hanya
memindahkan kendali tanpa menyediakan cara kembali; prosedur
melengkapinya dengan mekanisme pengembalian.

\begin{definisi}
\term{Prosedur} (disebut juga \term{subrutin}) adalah sekelompok
instruksi bernama yang dapat dipanggil dari bagian lain program,
menerima parameter, dan dapat mengembalikan nilai kepada
pemanggilnya.
\end{definisi}

Pemakaian prosedur tidak gratis. Setiap pemanggilan menimbulkan
\term{biaya pemanggilan prosedur}: mendorong parameter ke stack,
menjalankan \asm{CALL}, membangun lalu membongkar bingkai, dan
menjalankan \asm{RET}. Pada 8086, satu \asm{PUSH} memerlukan sekitar
\dec{11} sampai \dec{15} siklus, \asm{CALL} sekitar \dec{19} siklus,
dan \asm{RET} sekitar \dec{16} siklus. Untuk prosedur yang sangat
sederhana, biaya ini dapat melebihi biaya pekerjaannya sendiri.
Karena itu prosedur dipakai bukan terutama untuk mengejar kecepatan,
melainkan untuk kejelasan dan kemudahan pemeliharaan.

\subsection{Instruksi \asm{CALL} dan \asm{RET}}

\asm{CALL} mengerjakan dua hal sekaligus: menyimpan \term{alamat
kembali}, yaitu nilai \reg{CS} dan \reg{IP} pada saat instruksi
\emph{sesudah} \asm{CALL}, lalu memuat alamat prosedur yang dipanggil
ke \reg{CS} dan \reg{IP}.

\begin{lstlisting}[language=MASM, numbers=none]
    CALL HitungLuas         ; simpan alamat kembali, lalu lompat
    MOV  luas, AX           ; instruksi ini dieksekusi setelah RET
\end{lstlisting}

Alamat kembali yang disimpan pada contoh di atas adalah alamat
instruksi \lstasm|MOV luas, AX|. Ketika prosedur selesai, \asm{RET}
mengambil alamat itu dari stack dan memuatnya kembali ke \reg{IP}.
Stack dengan demikian merupakan bagian dari mekanisme kendali
program, bukan sekadar tempat menyimpan data.

\subsection{Pemanggilan \eng{Near} dan \eng{Far}}

Pada \term{pemanggilan dekat} (\eng{near call}), prosedur berada
pada segmen kode yang sama, sehingga hanya \reg{IP} yang perlu
disimpan. Pada \term{pemanggilan jauh} (\eng{far call}), prosedur
berada pada segmen lain, sehingga \reg{CS} dan \reg{IP} sama-sama
disimpan. Perbedaannya dirangkum pada Tabel~\ref{tab:call-near-far}.

\begin{table}[H]
    \centering
    \caption{Perbandingan pemanggilan \eng{near} dan \eng{far} pada
    prosesor 8086.}
    \label{tab:call-near-far}
    \small
    \begin{tabular}{@{}p{3.3cm}p{5.0cm}p{5.0cm}@{}}
        \toprule
        \textbf{Aspek} & \textbf{\asm{CALL} \eng{near}} &
        \textbf{\asm{CALL} \eng{far}} \\
        \midrule
        Isi yang didorong ke stack & \reg{IP} saja (\dec{2} byte) &
        \reg{CS} lalu \reg{IP} (\dec{4} byte) \\
        \addlinespace
        Instruksi kembali & \asm{RET} (\code{C3h}) &
        \asm{RETF} (\code{CBh}) \\
        \addlinespace
        Bentuk penulisan & \lstasm|CALL HitungLuas| &
        \lstasm|CALL LUAR:HitungLuas| \\
        \addlinespace
        Pemakaian lazim & Prosedur pada satu berkas program &
        Prosedur pada pustaka atau \eng{overlay} lain \\
        \bottomrule
    \end{tabular}
\end{table}

Seluruh contoh dalam bab ini memakai pemanggilan \eng{near} karena
prosedur dan pemanggilnya berada pada satu segmen kode.

\begin{infobox}
\asm{CALL} tidak sama dengan \asm{JMP}. Keduanya memindahkan
kendali, tetapi \asm{JMP} tidak menyimpan alamat kembali, sehingga
\asm{RET} akan mengambil alamat yang salah dari stack.
\end{infobox}

Setelah memahami bagaimana kendali berpindah, pertanyaan berikutnya
adalah di mana alamat kembali itu disimpan dan bagaimana isinya
berubah. Jawabannya adalah stack, yang dibahas pada
Subbab~\ref{sec:operasi-stack}.
